\chapter{Unit 5: Pandas DataFrame Operations \& Matplotlib Visualizations}

\begin{formulabox}{Pandas Selection \& Indexing Semantics: \texttt{.loc} vs \texttt{.iloc}}
\begin{center}
\begin{tabularx}{0.96\textwidth}{l X X}
\toprule
\textbf{Feature} & \textbf{\texttt{.loc} (Label-based Indexing)} & \textbf{\texttt{.iloc} (Integer Position Indexing)} \\
\midrule
Index specification & Row/Column index labels (strings, dates) & 0-based integer coordinates ($0 \dots N-1$) \\
Slice endpoint & \textbf{Inclusive of endpoint}: \texttt{'a':'c'} includes \texttt{'c'} & \textbf{Exclusive of endpoint}: \texttt{0:3} yields indices $0, 1, 2$ \\
Boolean masking & Supported: \texttt{df.loc[df['A'] > 0]} & Requires numpy array/list: \texttt{df.iloc[mask.values]} \\
Missing keys & Raises \texttt{KeyError} if label missing & Raises \texttt{IndexError} if out of bounds \\
\bottomrule
\end{tabularx}
\end{center}
\end{formulabox}

\begin{formulabox}{Data Cleaning \& Transformation APIs}
\begin{itemize}[leftmargin=1.5em, itemsep=2pt]
    \item \textbf{Missing Values}:
    \begin{itemize}
        \item Detection: \texttt{df.isna().sum()} (returns missing count per column).
        \item Dropping: \texttt{df.dropna(how='any', subset=['col1', 'col2'])}.
        \item Imputation: \texttt{df['val'].fillna(df['val'].median(), inplace=False)}.
    \end{itemize}
    \item \textbf{Deduplication}: \texttt{df.drop\_duplicates(subset=['id'], keep='first')}.
    \item \textbf{Categorical Conversion}: \texttt{df['cat'] = df['cat'].astype('category')} (reduces memory consumption by $80\%$).
    \item \textbf{Groupby Aggregation}:
    \begin{equation*}
    \texttt{df.groupby('dept')['salary'].agg(['mean', 'std', 'count']).reset\_index()}
    \end{equation*}
    \item \textbf{Transform vs Agg}: \texttt{.agg()} reduces group to single value; \texttt{.transform()} returns a Series aligned with original DataFrame rows.
\end{itemize}
\end{formulabox}

\begin{formulabox}{Matplotlib Object-Oriented Hierarchy Cheat-Sheet}
\begin{itemize}[leftmargin=1.5em, itemsep=2pt]
    \item Canvas initialization: \texttt{fig, axes = plt.subplots(nrows=2, ncols=2, figsize=(10, 8))}.
    \item Line chart: \texttt{ax.plot(x, y, color='navy', linestyle='--', linewidth=1.5, label='Trace')}.
    \item Scatter plot: \texttt{ax.scatter(x, y, s=area, c=colors, alpha=0.7)}.
    \item Histogram: \texttt{ax.hist(data, bins=30, density=True, edgecolor='black')}.
    \item Bar chart: \texttt{ax.bar(categories, values, color='orange')}, \texttt{ax.barh()} for horizontal.
    \item Labels \& Aesthetics:
    \begin{equation*}
    \texttt{ax.set\_title("Title"), ax.set\_xlabel("X"), ax.set\_ylabel("Y"), ax.legend(), ax.grid(True)}
    \end{equation*}
    \item Layout optimization: \texttt{plt.tight\_layout()}, \texttt{plt.savefig("fig.png", dpi=300, bbox\_inches='tight')}.
\end{itemize}
\end{formulabox}

\begin{memorytip}
Never modify a DataFrame via chained assignment \texttt{df[df['A']>0]['B'] = 5}! This triggers the dreaded \texttt{SettingWithCopyWarning} because Pandas cannot guarantee whether the intermediate slice is a view or a copy! Always write: \texttt{df.loc[df['A'] > 0, 'B'] = 5}.
\end{memorytip}
